\documentclass[10pt,a4paper]{amsart}
\usepackage[margin=19mm]{geometry}
\usepackage{amsmath,amssymb,mathtools}
\usepackage[scr=rsfs,cal=euler]{mathalfa}
\usepackage{xfrac}
\usepackage[unicode,hidelinks,bookmarks=false]{hyperref}
\newcommand{\ie}{i.e.\ }
\newcommand{\cf}{cf.\ }
\newcommand{\resp}{respectively\ }
\newcommand{\Subsection}{\subsection}
\providecommand{\nobreakdash}{\nobreak\hbox{-}}
\setlength{\emergencystretch}{2em}
\multlinegap0pt
\usepackage{bm}
\usepackage{bbm}
\usepackage{dsfont}
\usepackage{xspace}
\usepackage{cancel}
\usepackage{mathtools}
\usepackage{amsthm}
\makeatletter
\@ifundefined{theorem}{\newtheorem{theorem}{Theorem}}{}
\@ifundefined{corollary}{\newtheorem{corollary}[theorem]{Corollary}}{}
\@ifundefined{proposition}{\newtheorem{proposition}[theorem]{Proposition}}{}
\makeatother
\DeclareMathOperator{\diam}{diam}
\newcommand{\Z}{\mathbb Z}
\newcommand{\bo}{\boldsymbol}
\newcommand{\mc}{\mathcal}
\title[Filtering and Statistical Properties of Unimodal Maps Pertur]{Filtering and Statistical Properties of Unimodal Maps Perturbed by Heteroscedastic Noises}
\author{Fabrizio Lillo, Stefano Marmi, Matteo Tanzi, Sandro Vaienti}
\date{Source: arXiv:2411.13939v3 (CC BY 4.0)}
\begin{document}
\maketitle

\noindent\emph{Source and licence.} Filtering and Statistical Properties of Unimodal Maps Perturbed by Heteroscedastic Noises. Version arXiv:2411.13939v3 (CC BY 4.0), distributed under the Creative Commons Attribution 4.0 International licence, as recorded in the arXiv licence metadata for this exact version and shown on the abstract page https://arxiv.org/abs/2411.13939v3. Full rights notice: licenses/arxiv-2411.13939-CC-BY-4.0.txt.\par\smallskip

\noindent\textit{Editorial scope.} Counted unit: the Proposition whose source label is \texttt{None} in the article (source0.tex lines 597--599, source label \texttt{None}), delivered together with the article's own complete original proof of it (source0.tex lines 600--622, one \texttt{proof} environment of 23 lines). No mathematics is written, repaired or re-derived: statement and proof are the article's own text line for line. Required context retained uncounted: Eq:Evolution (lines 460--462); Eq:CocycleDef (lines 464--466); Cor:Contraction (lines 563--569). Every definition, display and label of those lines is retained unchanged. Omitted from this package and not redistributed: 1--459, 463--463, 467--562, 570--596, 623--1222. Nothing inside a retained excerpt is omitted or altered: every retained body is delivered exactly as the article's lines are written, so no omission receipt of any kind accompanies this package. Citations: the retained excerpts carry no citation command at all, so no bibliography entry of the article is transcribed and no reference is redistributed. Reference adaptation: none was needed and none was made; every reference inside the delivered unit resolves inside it, so no unresolved reference or citation is produced. Printed slips kept literal: no printed word, symbol, hypothesis or number of the counted statement or of its proof was altered. Layout and class: the article's own class is \texttt{article} with packages including amsmath, amssymb, amsthm, bbm, amsfonts, enumitem, tikz, comment, inputenc, isodate, caption, subcaption, multirow, fancyvrb; the wrapper is the lane's pinned \texttt{amsart} editorial wrapper at 10pt on A4, which restates the theorem-like environments the retained text needs and adds the article's own macro definitions that the retained text needs (\texttt{\textbackslash Z}, \texttt{\textbackslash bo}, \texttt{\textbackslash diam}, \texttt{\textbackslash mc}), with extra wrapper packages (bm, bbm, dsfont, xspace, cancel, mathtools). The delivered numbering is the unit's own. Assets: no retained span contains a figure environment, a graphics-inclusion command, a PDF-page inclusion, a table or a TikZ picture; no image file is copied into this package, the package declares no asset dependency and no image file is redistributed with it. Licence: version arXiv:2411.13939v3 of this article is distributed under the Creative Commons Attribution 4.0 International licence, as recorded in the arXiv licence metadata for this exact version and shown on the abstract page https://arxiv.org/abs/2411.13939v3. A full rights notice is delivered at licenses/arxiv-2411.13939-CC-BY-4.0.txt. Characters: every retained excerpt is pure ASCII, so the delivered engine, XeLaTeX with its default Latin Modern text font, reports no missing character.
\begin{equation}\label{Eq:Evolution}
(\mc P_z\nu)(\phi)= \frac{\int _I\phi(x)g(z,x)d\nu^+(x)}{\int _Ig(z,x)d\nu^+(x)}\quad\quad \nu^+(\phi)=\int K(x,\phi)d\nu(x).
\end{equation}

\begin{equation}\label{Eq:CocycleDef}
\mc P((\bo x,\bo z), \nu)=(\sigma(\bo x,\bo z),\mc P_{(\bo x,\bo z)}\nu).
\end{equation}

\begin{corollary}\label{Cor:Contraction}
If $\mc P$ is a linear transformation with $\mc P\mc V_0(J)\subset \mc V_c$ with $c\in(0,1]$, then
\[
\Theta_0(\mc P\phi,\mc P\psi)\le \lambda\Theta_0(\phi,\psi)\quad\quad\forall\phi,\psi\in\mc V_0
\]
with $\lambda=1-e^{-\diam_{\mc V_0(\mc V_c)}}\in[0,1)$.
\end{corollary}

\begin{proposition}
If the Main Assumption holds, there is a unique absolutely continuous equivariant probability measure for the cocycle in \eqref{Eq:CocycleDef}.
\end{proposition}

\begin{proof}
 With the above definitions, $\mc P_{(\bo x, \bo z)}$ sends nonnegative functions to nonnegative functions, i.e. it maps $\mc V_0(J_{z_0}) $ to $\mc V_0(J_{z_1})$.

Let $z_0\in I'$ and $z_1\in F_{z_0}'$ so that $J_{z_{1}}\subset F_{z_0}$, where $I'$, $F_{z_0}'$, and $F_{z_0}$ are as in the Main Assumption. Item 1) of the main assumption implies that given $\nu$ with density $\psi\in \mc V_0(J_{z_0})$, $\nu^+$ defined as in \eqref{Eq:Evolution} has density $\psi^+(y)=\int p(x,y)\psi(x)dx$ and  on $F_{z_{0}}$, $\psi^+>c$.  Since by assumption 2) we have that $g(z_{1},\cdot)>c$, for every $x,y\in J_{z_1}$, and by assumption 3) $g,\,\psi^+<c^{-1}$
\[
c^4<\frac{g(z_{1},x)\psi^+(x)}{g(z_{1},y)\psi^+(y)} < c^{-4}
\]
which means
\[
g(z_{1},\cdot)\psi^+(\cdot)\in \mc V_{c^4}(J_{z_1}).
\]
Therefore the linear application $\psi\mapsto g(z_{1},\cdot)\psi^+(\cdot)$ sends $\mc V_0(J_{z_0})$ to $\mc V_{c^4}(J_{z_{1}})$, by Corollary \ref{Cor:Contraction} it's a contraction with respect to the Hilbert metric $\Theta_0$. Since when $\nu$ has density $\psi$, by definition, $\mc P_{(\bo x,\bo z)}\nu$ has density
\[
\frac{g(z_{1},\cdot)\psi^+(\cdot)}{\int_Ig(z_{1},y)\psi^+(y)dy}
\]
and the Hilbert metric is projective, then also $\mc P_{(\bo x,\bo z )}$ from $\mc V_0(J_{z_0})$ to  $\mc V_{c^4}(J_{z_1})\subset\mc V_{0}(J_{z_1})$ is a contraction with rate $\lambda\in(0,1)$ independent of $z_0\in I'$.

Now consider the event $E:=\{Z_n\in I' \mbox{ and } J_{Z_{n+1}}\subset F_{Z_n} \}$. The main assumption implies that $E$ has positive measure and, by ergodicity, this event will occur almost surely for infinitely many $n\in \Z$ and with positive frequency.  This implies that for $\bo \mu$-a.e. sequence $(\bo x, \bo z)$ the set
\[
\bigcap_{n\ge 1} \mc P^n_{(\sigma^{-n}\bo x,\sigma^{-n}\bo z)} \mc V_0(J_{z_{-n}})\
\]
contains a unique direction: the intersection is non-empty because intersection of compact (with respect to the weak topology) projective sets; the direction is unique, because the diameter of $\mc P^n_{(\sigma^{-n}\bo x,\sigma^{-n}\bo z)} \mc V_0(J_{z_{-n}})$ tends to zero for a.e. $(\bo x, \bo z)$. The probability density in this unique direction is the equivariant measure in the fibre sitting on $(\bo x, \bo z)$ (the equivariance is immediate by construction).
\end{proof}

\end{document}
